\section{The permanent seventh minimum and the bound forty-three}

\begin{finding}{R-10}{the exact permanent seventh, eighth and ninth minima}
Over every algebraically closed field of characteristic zero,
\[
 \mu_7(U^{\mathrm{per}}_{4,2})=14,\qquad
 \mu_8(U^{\mathrm{per}}_{4,2})=15,\qquad
 \mu_9(U^{\mathrm{per}}_{4,2})=16,
 \qquad \BR(\mathcal P_4)\geq43.
\]
The same lower bound holds for the usual complex symmetric border
rank of the four-by-four permanent polynomial.
\end{finding}
\begin{proof}
P-01 reduces the seventh Grassmannian minimum to coordinate
seven-planes among the $36$ complementary permanent minors.
There are $\binom{36}{7}=8\,347\,680$ such subsets. For a subset
$S$, let $m(S)$ be the number of matrix-entry variables appearing
in its seven rectangular supports. The ordinary support flattening
has rank $m(S)$. Indeed, for each appearing variable $x$, choose
one rectangle $s$ containing it. The column indexed by $s$ and
the diagonally opposite variable has a single unit in row $x$;
these chosen columns are distinct for distinct $x$. They form
a unit minor of order $m(S)$. Thus $m(S)\geq14$ gives the required
lower bound directly.

Row and column permutations, together with transposition, act on
the $36$ rectangles and preserve the permanent tensor up to
invertible changes of bases. The resulting group has order
$24^2\cdot2=1152$. The remaining $1\,543\,608$ subsets, those with
$m(S)<14$, form exactly $1517$ orbits. The certificate proves this
finite cover rather than relying on a classification assertion:
it constructs all $1152$ actions, verifies that the $1517$ listed
orbits are pairwise disjoint and have the recorded sizes, and
independently enumerates all $\binom{36}{7}$ subsets to check that
every subset of support less than fourteen belongs to their union.

The lower certificates for these $1517$ representatives are:
\begin{center}
\begin{tabular}{p{0.39\linewidth}|r|p{0.40\linewidth}}
Certificate type&Orbits&Lower certificate\\\hline
Koszul map on the seven-dimensional factor&1503&Rank $>13\binom6p$\\
Koszul map on the variable-support factor&3&Rank $>13\cdot495$\\
P-13&5&Support-twelve prolongation exclusions\\
P-15&2&Support-nine prolongation exclusions\\
P-16&3&Support-eleven prolongation exclusions\\
P-19&1&The last higher-prolongation exclusion
\end{tabular}
\end{center}
For the three support-factor Koszul maps, $m=13$, the exterior
degree is $p=4$, and the rank of the map on a pure tensor is
$\binom{12}{4}=495$. Their certified ranks are respectively
$6457$, $6457$ and $6515$, all strictly greater than
$13\cdot495=6435$. Each rank is a sum of certified ranks of
distinct torus weight blocks. For the first-factor maps, each
listed minor has order strictly greater than thirteen times
the corresponding pure-tensor cost. Determinantal closedness
therefore gives border rank at least fourteen in all these cases.
The eleven remaining representatives are matched, as exact
sorted subsets of rectangle indices, with the separately proved
P-13, P-15, P-16 and P-19 statements.

The file \path{per7_global_certificate.json} records the full
cover, weights and minor indices. The independent verifier
\path{verify_per7_global.py} reconstructs the tensors and all
selected matrix entries, using only the audited arithmetic and
Koszul-entry routines in \path{verify_det7_global.py}.
It verifies $6041$ nonzero minors modulo the prime $1000003$
(whose primality is checked by trial division), including
disjointness and the asserted weights of the blocks. A nonzero
residue proves that the corresponding integer determinant is
nonzero, hence gives the claimed rank over every field of
characteristic zero. This statement does not infer the same
rank in arbitrary positive characteristic from a modular
calculation. The exhaustive cover and the eleven separately
checked prolongation proofs establish $\mu_7\geq14$.

An anchored $3\times3$ shape with seven edges meeting every
row and column vertex has border rank $7+3+3+1=14$ by P-06.
It is a coordinate seven-plane in this residual space, so
$\mu_7\leq14$. R-07 now gives $\mu_j\geq j+7$ for $j\geq7$.
The eight- and nine-edge anchored shapes give equality for
$j=8,9$, again by P-06. Finally, the residual image has
dimension $36$, and R-01 gives
\[
 \BR(\mathcal P_4)\geq36-j+\mu_j=43\qquad(j=7,8,9).
\]
In characteristic zero the integral coefficient tensor differs
from the normalized polynomial tensor by a nonzero scalar.
The symmetric atom variety is contained in the CP atom variety,
so the CP lower bound also bounds symmetric border rank.
\end{proof}

\paragraph{Geometric dictionary.}
R-10 uses the same relative Grassmannian, torus degeneration,
exterior morphisms of locally free sheaves and proper
prolongation incidences as the preceding findings. For the
permanent, the finite symmetry group and exhaustive coordinate
cover replace the determinant's Borel order-ideal reduction;
no characteristic-zero Borel invariance of the permanent is
assumed. The minors are functions defined over $\Z$, and their
nonzero residues certify nonvanishing over $\Q$ and then over
every characteristic-zero extension. The maps and their
determinantal ideals, rather than their $K_0$ classes alone,
carry the rank information. Exact foundational locators and
the absence of an established archival cote for these modern
calculations are recorded in the separate dictionary.
